% !TEX root = ../main.tex

\chapter{大规模关系表格数据可扩展训练框架 RelX 的设计与实现}
\label{chap:relx-design}

\section{问题定义与设计目标}
\thesisplaceholder{界定大规模关系表格训练的输入、资源开销与可扩展性目标；分析数据规模增长时全量特征物化的内存开销，以及逐批准备对训练吞吐的影响。}

\section{RelX 总体架构}
\thesisplaceholder{绘制离线准备与在线训练的总体架构图，说明采样、特征准备和模型训练三个在线阶段的数据流及模块边界。}

\section{离线结构与统计准备}
\thesisplaceholder{说明模式、拓扑和必要统计量的准备过程，区分必须常驻的结构与按需获取的原始行及编码特征。}

\section{采样计划与按需特征读取}
\thesisplaceholder{定义采样计划的数据结构、节点请求去重和逆映射，说明如何从数据库读取所需行并恢复原批次的节点顺序。}

\section{异构字段编码与批次组装}
\thesisplaceholder{说明数值、类别、文本字段从原始行到模型输入的转换、组装进程的职责以及跨进程传递的对象。}

\section{分层缓存与特征复用}
\thesisplaceholder{描述原始行、已编码行、文本向量等缓存层的键、准入、失效和预算；说明不同负载下可能的收益与代价。}

\section{有界并行流水线}
\thesisplaceholder{说明采样、组装和训练的生产者—消费者关系、有界队列、在途窗口、执行顺序及异常处理，并区分批内依赖与跨批重叠。}

\section{语义一致性与适用范围}
\thesisplaceholder{列明任务划分、时间可见性、邻居采样、节点标识、字段类型和编码规则等不变性约束；说明当前拓扑驻内存等实现边界。}

\section{本章小结}
\thesisplaceholder{归纳 RelX 在数据规模增长条件下控制内存开销、减少重复准备和维持训练吞吐的设计，连接下一章的验证实验。}
